\scanpage{325}
\mspage{64}
% archive sheet number: 324
% ms. 1–3
\deleted{\unsure{La démonstration de}} Notons que \inserted{plus \unsure{génért}}, pour tout chemin $c\unsure{;:|:}\,a\longrightarrow b$, \inserted{de type $\tau$,} on définit un \uline{chemin opposé}
% ms. 4
\[c^{\unsure{\circ|o}}:b\longrightarrow a\qquad\text{de type }\tau^\circ,\]
% ms. 5
où
% ms. 6
\[\tau\longmapsto\tau^\circ:\quad\mathrm{Ty}\longrightarrow\mathrm{Ty}\]
% ms. 7–10
est l'\unsure{anti-involution} \deleted{\unsure{de}} du monoïde libre $\mathrm{Ty}$ qui échange les deux générateurs $\Delta_-,\Delta_+$, et qui \deleted{\unsure{sur}} est donc donné sur l'élément général par
% ms. 11
\[(\Delta_{\varepsilon_1}^{\deleted{\unsure{\varepsilon_1}}}\,\Delta_{\varepsilon_2}^{\deleted{\unsure{\varepsilon_2}}}\,\cdots\,\Delta_{\varepsilon_n})^\circ
=\Delta_{-\varepsilon_n}\,\Delta_{-\varepsilon_{n-1}}\cdots\Delta_{-\varepsilon_{\unsure{n-|n|1}}}.\]
% ms. 12–15
Par récurrence \uline{sur la longueur} de $c$, en \unsure{utilisant} \deleted{\unsure{enfin}} le résultat précédent \deleted{que} que $c(u)$, $c(u)^\circ$ sont inverses l'un de l'autre dans $\Pi(X)$, on trouve
% ms. 12–15 left
% slanted
\margin{inutile, \unsure{résulte} directement de la formule $(*)$ page suivante, \unsure{en} décomposant $c$ en \unsure{composé|suite} de chemins ``\unsure{atomiques|élémentaires}'' de type $\Delta_+$ ou $\Delta_-$.}
% ms. 16–17
\[\left.\begin{aligned}c^\circ c&=\mathrm{id}_a\\cc^\circ&=\mathrm{id}_b\end{aligned}\right)\quad\text{dans }\Pi(X).\]

% ms. 18–19
\uline{Exemple} Si
% ms. 18–19
\[\begin{aligned}\tau&=\Delta_+\,\Delta_-\,\Delta_+^2\,\Delta_{\unsure{-|+}}\\
\tau'&=\Delta_+\,\Delta_-^2\,\Delta_+\,\Delta_-\end{aligned}\]
% ms. 20
Et si on a un chemin
% ms. 21
\diagram{\begin{tikzcd}[ampersand replacement=\&,column sep=small]
c:\quad a_0\arrow[r,"u_1"] \& a_1 \& a_2\arrow[l,"u_2"']\arrow[r,"u_3"] \& a_3\arrow[r,"u_4"] \& a_4 \& a_5\arrow[l,"u_5"']
\end{tikzcd}}
% ms. 22
de type $\tau$, on en déduit \deleted{\unsure{des}} le chemin
% ms. 23
\diagram{\begin{tikzcd}[ampersand replacement=\&,column sep=small]
\underline{c}^\circ:\quad a_5\arrow[r,"u_5"] \& a_4 \& a_3\arrow[l,"u_4"'] \& a_2\arrow[l,"u_3"']\arrow[r,"u_2"] \& a_1 \& a_0\arrow[l,"u_1"']
\end{tikzcd}}
% ms. 24
de type $\tau'$. L'opération $c\longmapsto c^\circ$ satisfait en
